\section{Impactos sociales y ambientales del proyecto}

\subsection{Impacto social y economico}

El proyecto presenta principalmente un impacto social, dado que se trata de un videojuego de tipo MMO RPG. La masividad 
del mismo permitiría a usuarios y empresas promocionar productos y servicios. No solo creando mundos por ellos mismos y 
permitiendo que los usuarios jueguen en ellos, sino también con la colaboración de estos agentes, a través del equipo 
desarrollador.
\vspace{0.25cm}

Esta estrategia es utilizada por videojuegos como \textit{Fortnite}. En este caso, la empresa creadora del juego (\textit{Epic Games}) 
realiza colaboraciones con empresas y personalidades del momento, promocionando \textit{skins} o realizando eventos dentro 
del juego. De esa forma se potencian ambos actores de manera sinérgica: \textit{Epic Games} promocionando su juego y los 
colaboradores generando marca a través de los seguidores del juego. Esta estrategia podría ser utilizada, del mismo modo, 
en \textit{Feed the Realm}, realizando colaboraciones en las que se crean mundos personalizados y cosméticos dentro del 
juego.
\vspace{0.25cm}

Por otro lado, al tratarse de un videojuego, si se realiza una correcta estrategia de marketing y el proyecto se mantiene 
añadiendo más y mejores funcionalidades, el mismo podría llamar la atención de \textit{streamers} o \textit{gamers} que 
juegan o crean sus mundos en plataformas de streaming como \textit{YouTube}, \textit{Twitch} o \textit{Kick}. De esa forma 
se podrán atraer interesados que quieran probar el juego y/o quieran crear sus propios mundos.
\vspace{0.25cm}

Afortunadamente, el impacto social no es un fin en sí mismo, sino que es el medio para conseguir otros fines. Dada la masividad 
del proyecto, el mismo permitirá crear una comunidad. La comunidad no es solo personas usando el producto, sino también 
la posibilidad de mantener vivo el proyecto. Esta dará \textbf{\textit{feedback} inmediato}, lo que permitirá mejorar 
las funcionalidades existentes y agregar una mayor cantidad de \textit{features} de las que ya se dispone. Este ciclo de 
retroalimentación positiva permitirá que las personas se interesen por el videojuego, seguido de una necesidad manifiesta 
por parte del usuario que demandará más y mejores funcionalidades, lo que finalmente atraerá a más personas a probar el 
producto.
\vspace{0.25cm}

Este impacto social trae aparejado consigo un impacto económico, dado que el beneficio del proyecto viene dado por jugadores 
que compren cosméticos o creadores de mundos que se suscriban para desplegar sus mundos. Si el proyecto alcanza un impacto 
social, mediante correctas estrategias de marketing y un mantenimiento sostenido a lo largo del tiempo, el impacto económico 
será alcanzado.
\vspace{0.25cm}

Es importante destacar que estas no son las únicas estrategias para alcanzar el deseado impacto económico. Las colaboraciones 
\textbf{\textit{Business to Business} (B2B)} que el equipo realice con personalidades y/o diferentes empresas pueden venir 
con contratos en los que se comprometa a ofrecer \textit{features} personalizadas, servidores dedicados al alojamiento de 
sus mundos y/o la realización de los mundos por ellos. Todo ello, a cambio de un beneficio económico, y ambas partes recibirían 
notoriedad.
\vspace{0.25cm}

Sin embargo, no todo el impacto social debe traducirse inmediatamente en un beneficio económico. Dada la libertad de 
creación de mundos, entidades gubernamentales y/o educativas podrían utilizar esa libertad para crear mundos con el fin 
de fomentar en la sociedad un tema de interés del momento, o bien las escuelas podrían crear mundos para facilitar la 
comprensión de un tema en el aula, a través del videojuego. El equipo barajó, por ejemplo, la posibilidad de crear un 
mundo ambientado en la Semana del 25 de Mayo, pero también se podrían recrear batallas históricas o diferentes hitos de 
la humanidad.
\vspace{0.25cm}

Cabe detenerse también en una cuestión coyuntural: jugar un videojuego es una actividad de ocio y parece no tener un impacto 
notorio aparente más que distraerse un rato o pasar un buen momento. Sin embargo, luego de la \textbf{pandemia de COVID-19}, 
la sociedad en general, a lo largo de todo el mundo, ha preferido mantenerse más aislada y socializando cada vez menos. 
Jugar un videojuego online no es solamente evadir un mal momento o divertirse, sino compartir momentos con otros. También 
existe la posibilidad de conectar a diferentes personas a través de los chats, con lo cual se podría fomentar el encuentro 
de amigos o grupos de amigos de forma presencial, desalentando la falta de sociabilización del último tiempo.

\subsection{Impacto Ambiental}
